% ============================================================================
%  estilo.tex — identidade visual, tipografia e componentes do documento
%
%  Tipografia: Roboto (texto e titulos) + Roboto Mono (codigo), conforme
%  solicitado. Ambas as familias vem do pacote `roboto` (Type1/OTF), portanto
%  o documento compila com pdfLaTeX puro — sem fontspec e sem XeLaTeX.
% ============================================================================

% --- Codificacao e idioma ---------------------------------------------------
\usepackage[T1]{fontenc}
\usepackage[utf8]{inputenc}
\usepackage[brazil]{babel}
\usepackage{csquotes}

% --- Tipografia -------------------------------------------------------------
\usepackage[sfdefault]{roboto}            % Roboto como familia padrao
\usepackage[scaled=0.90]{roboto-mono}     % Roboto Mono para codigo
\usepackage{microtype}
\usepackage{textcomp}

% --- Pagina -----------------------------------------------------------------
\usepackage[
  a4paper,
  top=2.6cm, bottom=2.5cm, left=2.6cm, right=2.4cm,
  headheight=15pt, headsep=18pt, footskip=26pt
]{geometry}
\usepackage{pdflscape}
\usepackage{setspace}
\onehalfspacing

% --- Cores (paleta validada do projeto — src/ui/theme.py) -------------------
\usepackage{xcolor}
\definecolor{pmxBlue}      {HTML}{2A78D6}  % primaria / documentado
\definecolor{pmxBlueDark}  {HTML}{184F95}  % primaria escura
\definecolor{pmxBlueLight} {HTML}{86B6EF}
\definecolor{pmxBlueWash}  {HTML}{EAF2FC}
\definecolor{pmxOrange}    {HTML}{EB6834}  % sem documento / atencao
\definecolor{pmxOrangeWash}{HTML}{FDEDE7}
\definecolor{pmxGreen}     {HTML}{1BAF7A}
\definecolor{pmxGreenWash} {HTML}{E6F7F1}
\definecolor{pmxYellow}    {HTML}{EDA100}
\definecolor{pmxYellowWash}{HTML}{FDF3DC}
\definecolor{pmxPink}      {HTML}{E87BA4}
\definecolor{pmxRed}       {HTML}{D03B3B}
\definecolor{pmxRedWash}   {HTML}{FBEAEA}
\definecolor{pmxInk}       {HTML}{0B0B0B}  % texto primario
\definecolor{pmxInkSoft}   {HTML}{52514E}  % texto secundario
\definecolor{pmxMuted}     {HTML}{898781}
\definecolor{pmxGrid}      {HTML}{E1E0D9}
\definecolor{pmxAxis}      {HTML}{C3C2B7}
\definecolor{pmxSurface}   {HTML}{FCFCFB}
\definecolor{pmxPaper}     {HTML}{F5F5F1}

% --- Tabelas ----------------------------------------------------------------
\usepackage{booktabs}
\usepackage{tabularx}
\usepackage{longtable}
\usepackage{array}
\usepackage{multirow}
\usepackage{makecell}
\usepackage{colortbl}
\usepackage{ltablex}          % tabularx com quebra de pagina
\keepXColumns
\renewcommand{\arraystretch}{1.22}
\arrayrulecolor{pmxAxis}
\setlength{\aboverulesep}{0pt}
\setlength{\belowrulesep}{0pt}
\setlength{\extrarowheight}{2pt}

\newcolumntype{L}[1]{>{\raggedright\arraybackslash}p{#1}}
\newcolumntype{C}[1]{>{\centering\arraybackslash}p{#1}}
\newcolumntype{R}[1]{>{\raggedleft\arraybackslash}p{#1}}
\newcolumntype{Y}{>{\raggedright\arraybackslash}X}

% Colunas X sem justificacao: as celulas misturam texto e codigo monoespacado,
% que nao hifeniza — justificar produziria espacos internos muito largos.
% \arraybackslash preserva o \\ como separador de linha na ultima coluna.
\renewcommand{\tabularxcolumn}[1]{>{\raggedright\arraybackslash}p{#1}}

\newcommand{\thd}[1]{\textcolor{white}{\textbf{#1}}}
\newcommand{\headrow}{\rowcolor{pmxBlueDark}}
\newcommand{\zebra}{\rowcolor{pmxPaper}}

% --- Listas -----------------------------------------------------------------
\usepackage{enumitem}
\setlist[itemize]{leftmargin=1.35em, itemsep=2pt, topsep=4pt, parsep=0pt,
                  label=\textcolor{pmxBlue}{\raisebox{0.18ex}{\rule{0.40em}{0.40em}}}}
\setlist[enumerate]{leftmargin=1.7em, itemsep=2pt, topsep=4pt, parsep=0pt,
                    font=\bfseries\color{pmxBlueDark}}
\setlist[description]{leftmargin=1.4em, style=nextline,
                      font=\bfseries\color{pmxBlueDark}}

% --- Matematica e unidades --------------------------------------------------
\usepackage{amsmath}
\usepackage{amssymb}
\usepackage{siunitx}
\sisetup{output-decimal-marker={,}, group-separator={.}, group-minimum-digits=4,
         per-mode=symbol, detect-weight=true, detect-family=true}
% Unidades nao-SI presentes no conjunto de dados original (colunas redundantes)
\DeclareSIUnit{\inch}{in}
\DeclareSIUnit{\degF}{\text{\textdegree}F}
\DeclareSIUnit{\gforce}{g}

% --- Figuras e legendas -----------------------------------------------------
\usepackage{graphicx}
\usepackage{float}
\usepackage{adjustbox}
\usepackage[font=small, labelfont={bf,color=pmxBlueDark}, textfont=normalfont,
            labelsep=endash, justification=raggedright, singlelinecheck=false,
            skip=8pt]{caption}
\usepackage{subcaption}
\captionsetup[table]{position=top, skip=6pt}

% Ajusta a figura a largura do texto apenas quando ela excede a caixa.
\newcommand{\fitwidth}[1]{\adjustbox{max width=\linewidth, center}{#1}}
\newcommand{\fitpage}[1]{\adjustbox{max width=\linewidth, max height=0.82\textheight, center}{#1}}

% --- Blocos de destaque (tcolorbox) ----------------------------------------
\usepackage[most]{tcolorbox}
\tcbuselibrary{skins,breakable}

\tcbset{
  pmxbase/.style={
    enhanced, breakable, sharp corners,
    boxrule=0pt, leftrule=2.6pt, arc=0pt,
    left=10pt, right=10pt, top=7pt, bottom=7pt,
    fonttitle=\bfseries\small,
    before skip=10pt, after skip=10pt
  }
}

\newtcolorbox{nota}[1][]{pmxbase, colback=pmxBlueWash, colframe=pmxBlue,
  coltitle=pmxBlueDark, title={NOTA}, #1}
\newtcolorbox{atencao}[1][]{pmxbase, colback=pmxOrangeWash, colframe=pmxOrange,
  coltitle=pmxOrange, title={ATENÇÃO}, #1}
\newtcolorbox{decisao}[1][]{pmxbase, colback=pmxGreenWash, colframe=pmxGreen,
  coltitle=pmxGreen, title={DECISÃO DE PROJETO}, #1}
\newtcolorbox{limitacao}[1][]{pmxbase, colback=pmxRedWash, colframe=pmxRed,
  coltitle=pmxRed, title={LIMITAÇÃO CONHECIDA}, #1}
\newtcolorbox{destaque}[1][]{pmxbase, colback=pmxPaper, colframe=pmxInkSoft, #1}
\newtcolorbox{especificacao}[1]{pmxbase, colback=pmxSurface, colframe=pmxBlueDark,
  coltitle=pmxBlueDark, title={#1}}

% --- Codigo-fonte -----------------------------------------------------------
\usepackage{listings}

\definecolor{lstBg}     {HTML}{FAFAF8}
\definecolor{lstFrame}  {HTML}{E1E0D9}
\definecolor{lstKeyword}{HTML}{184F95}
\definecolor{lstString} {HTML}{0F7B57}
\definecolor{lstComment}{HTML}{898781}
\definecolor{lstEmph}   {HTML}{C2410C}

\lstdefinestyle{pmx}{
  backgroundcolor=\color{lstBg},
  basicstyle=\ttfamily\scriptsize\color{pmxInk},
  keywordstyle=\color{lstKeyword}\bfseries,
  keywordstyle=[2]\color{lstEmph}\bfseries,
  keywordstyle=[3]\color{pmxBlue},
  stringstyle=\color{lstString},
  commentstyle=\color{lstComment}\itshape,
  numberstyle=\ttfamily\tiny\color{pmxMuted},
  identifierstyle=\color{pmxInk},
  numbers=left, numbersep=9pt, stepnumber=1,
  frame=single, framesep=6pt, rulecolor=\color{lstFrame},
  xleftmargin=17pt, xrightmargin=2pt,
  breaklines=true, breakatwhitespace=false, breakindent=0pt,
  postbreak=\mbox{\textcolor{pmxBlue}{$\hookrightarrow$}\space},
  showstringspaces=false, tabsize=4,
  columns=fullflexible, keepspaces=true, upquote=true,
  aboveskip=10pt, belowskip=10pt,
  literate=
    {á}{{\'a}}1 {à}{{\`a}}1 {ã}{{\~a}}1 {â}{{\^a}}1
    {é}{{\'e}}1 {ê}{{\^e}}1 {í}{{\'i}}1
    {ó}{{\'o}}1 {õ}{{\~o}}1 {ô}{{\^o}}1
    {ú}{{\'u}}1 {ü}{{\"u}}1 {ç}{{\c c}}1
    {Á}{{\'A}}1 {À}{{\`A}}1 {Ã}{{\~A}}1 {Â}{{\^A}}1
    {É}{{\'E}}1 {Ê}{{\^E}}1 {Í}{{\'I}}1
    {Ó}{{\'O}}1 {Õ}{{\~O}}1 {Ô}{{\^O}}1
    {Ú}{{\'U}}1 {Ç}{{\c C}}1
    {°}{{\textdegree}}1 {º}{{\textordmasculine}}1
    {µ}{{$\mu$}}1 {×}{{$\times$}}1 {≈}{{$\approx$}}1
    {≥}{{$\geq$}}1 {≤}{{$\leq$}}1 {≠}{{$\neq$}}1
    {→}{{$\rightarrow$}}1 {←}{{$\leftarrow$}}1 {⇔}{{$\Leftrightarrow$}}1
    {—}{{---}}3 {–}{{--}}2 {…}{{\ldots}}3
    {“}{{\textquotedblleft}}1 {”}{{\textquotedblright}}1
    {‘}{{\textquoteleft}}1 {’}{{\textquoteright}}1
    {ª}{{\textordfeminine}}1 {§}{{\S}}1
}
\lstset{style=pmx}

\lstdefinelanguage{yamlpmx}{
  morekeywords={true,false,null,families,rules,pattern,family,is_fault,display,
            title,services,image,build,command,environment,volumes,ports,
            depends_on,healthcheck,condition,runs-on,steps,uses,needs,jobs,
            name,on,with,env,test,interval,timeout,retries},
  sensitive=true,
  morecomment=[l]{\#},
  morestring=[b]",
  morestring=[b]'
}
\lstdefinelanguage{mermaid}{
  morekeywords={flowchart,graph,sequenceDiagram,erDiagram,classDiagram,gantt,
            subgraph,end,participant,actor,alt,else,loop,par,opt,note,activate,
            deactivate,direction,TB,LR,RL,BT,Person,Person_Ext,System,System_Ext,
            SystemDb,Container,ContainerDb,ContainerQueue,Component,ComponentDb,
            Rel,Rel_D,Rel_U,Rel_L,Rel_R,System_Boundary,Container_Boundary,
            Enterprise_Boundary,Boundary,Deployment_Node,UpdateLayoutConfig,
            UpdateRelStyle,classDef,class,style,state,C4Context,C4Container,
            C4Component,C4Deployment,stateDiagram},
  sensitive=true,
  morecomment=[l]{\%\%},
  morestring=[b]"
}
\lstdefinelanguage{jsonpmx}{
  morekeywords={true,false,null},
  keywordstyle=\color{lstEmph}\bfseries,
  sensitive=true,
  morestring=[b]"
}
\lstdefinelanguage{tspmx}[]{Java}{
  morekeywords={const,let,export,default,import,from,as,interface,type,await,
                async,function,return,string,number,boolean,readonly},
  sensitive=true,
  morecomment=[l]{//},
  morecomment=[s]{/*}{*/},
  morestring=[b]",
  morestring=[b]'
}

% Codigo em linha.
%
% Uso obrigatorio de \texttt e \textcolor (e nao de {\ttfamily\color{...}}):
% os comandos \text... emitem \leavevmode antes de trocar a fonte, o que faz o
% paragrafo comecar na fonte do texto corrente. Sem isso, uma celula de tabela
% que inicia com codigo teria a primeira linha de base calculada com as metricas
% da fonte monoespacada e ficaria visivelmente desalinhada das demais colunas.
\newcommand{\cd}[1]{\texttt{\textcolor{pmxBlueDark}{#1}}}
\newcommand{\arq}[1]{\texttt{\textcolor{pmxInkSoft}{#1}}}

% --- Titulos ----------------------------------------------------------------
\usepackage{titlesec}
\usepackage{titletoc}

\titleformat{\chapter}[display]
  {\normalfont}
  {\raggedleft\fontsize{62}{62}\selectfont\color{pmxBlueLight}\textbf{\thechapter}}
  {-10pt}
  {\raggedright\fontsize{23}{27}\selectfont\bfseries\color{pmxBlueDark}}
  [\vspace{5pt}{\color{pmxBlue}\rule{\linewidth}{2.2pt}}]
\titlespacing*{\chapter}{0pt}{-30pt}{26pt}

\titleformat{name=\chapter,numberless}[display]
  {\normalfont}
  {}
  {0pt}
  {\raggedright\fontsize{23}{27}\selectfont\bfseries\color{pmxBlueDark}}
  [\vspace{5pt}{\color{pmxBlue}\rule{\linewidth}{2.2pt}}]

\titleformat{\section}
  {\normalfont\Large\bfseries\color{pmxBlueDark}}{\thesection}{0.7em}{}
\titleformat{\subsection}
  {\normalfont\large\bfseries\color{pmxInk}}{\thesubsection}{0.7em}{}
\titleformat{\subsubsection}
  {\normalfont\normalsize\bfseries\color{pmxInkSoft}}{\thesubsubsection}{0.7em}{}
\titleformat{\paragraph}[runin]
  {\normalfont\normalsize\bfseries\color{pmxBlueDark}}{}{0pt}{}[\;---\;]

\titlespacing*{\section}{0pt}{20pt}{8pt}
\titlespacing*{\subsection}{0pt}{15pt}{6pt}
\titlespacing*{\subsubsection}{0pt}{12pt}{5pt}
\titlespacing*{\paragraph}{0pt}{9pt}{0.4em}

\setcounter{secnumdepth}{3}
\setcounter{tocdepth}{2}

% --- Sumario ----------------------------------------------------------------
\titlecontents{chapter}[0em]
  {\addvspace{9pt}\bfseries\color{pmxBlueDark}}
  {\contentslabel[\thecontentslabel]{1.5em}}
  {}
  {\hfill\contentspage}
\titlecontents{section}[1.7em]
  {\addvspace{1.5pt}}
  {\contentslabel{2.2em}}
  {}
  {\titlerule*[0.7pc]{\color{pmxAxis}.}\contentspage}
\titlecontents{subsection}[4.0em]
  {\addvspace{0.5pt}\small\color{pmxInkSoft}}
  {\contentslabel{2.8em}}
  {}
  {\titlerule*[0.7pc]{\color{pmxGrid}.}\contentspage}

% --- Cabecalho e rodape -----------------------------------------------------
\usepackage{fancyhdr}
\usepackage{lastpage}

\newcommand{\docid}{PMX-DOC-001}
\newcommand{\docversao}{1.0}

\renewcommand{\headrule}{{\color{pmxGrid}\hrule height \headrulewidth}}
\renewcommand{\footrule}{{\color{pmxGrid}\hrule height \footrulewidth}}

\fancypagestyle{pmxmain}{%
  \fancyhf{}
  \fancyhead[L]{\footnotesize\color{pmxInkSoft}\nouppercase{\leftmark}}
  \fancyhead[R]{\footnotesize\color{pmxMuted}Manutenção Prescritiva}
  \fancyfoot[L]{\footnotesize\color{pmxMuted}\docid{} \textperiodcentered{} v\docversao}
  \fancyfoot[R]{\footnotesize\color{pmxInkSoft}\thepage\,/\,\pageref{LastPage}}
  \renewcommand{\headrulewidth}{0.5pt}
  \renewcommand{\footrulewidth}{0.5pt}
}
% Pre-textual: numeracao romana, sem total de paginas (o total do
% \pageref{LastPage} refere-se a numeracao arabe do corpo do documento).
\fancypagestyle{pmxfront}{%
  \fancyhf{}
  \fancyhead[L]{\footnotesize\color{pmxInkSoft}\nouppercase{\leftmark}}
  \fancyhead[R]{\footnotesize\color{pmxMuted}Manutenção Prescritiva}
  \fancyfoot[L]{\footnotesize\color{pmxMuted}\docid{} \textperiodcentered{} v\docversao}
  \fancyfoot[R]{\footnotesize\color{pmxInkSoft}\thepage}
  \renewcommand{\headrulewidth}{0.5pt}
  \renewcommand{\footrulewidth}{0.5pt}
}
\fancypagestyle{plain}{%
  \fancyhf{}
  \fancyfoot[R]{\footnotesize\color{pmxInkSoft}\thepage}
  \renewcommand{\headrulewidth}{0pt}
  \renewcommand{\footrulewidth}{0pt}
}
\pagestyle{pmxfront}
\renewcommand{\chaptermark}[1]{\markboth{\thechapter\quad #1}{}}

% --- Referencias cruzadas e PDF --------------------------------------------
\usepackage[
  colorlinks=true,
  linkcolor=pmxBlueDark,
  citecolor=pmxBlue,
  urlcolor=pmxBlue,
  filecolor=pmxBlue,
  bookmarksnumbered=true,
  bookmarksopen=true,
  bookmarksopenlevel=1,
  pdfdisplaydoctitle=true
]{hyperref}
\usepackage{bookmark}
\usepackage{url}
\urlstyle{sf}

% --- Bibliografia (ABNT NBR 6023) ------------------------------------------
\usepackage[backend=biber, style=abnt, giveninits=true]{biblatex}
\addbibresource{referencias.bib}

% --- Semantica de referencia do projeto ------------------------------------
\newcommand{\req}[1]{\textbf{\textcolor{pmxBlueDark}{#1}}}
\newcommand{\adrr}[1]{\textbf{\textcolor{pmxGreen}{#1}}}
\newcommand{\ucr}[1]{\textbf{\textcolor{pmxBlue}{#1}}}
\newcommand{\yes}{\textcolor{pmxGreen}{\textbf{\checkmark}}}
\newcommand{\no}{\textcolor{pmxOrange}{\textbf{--}}}

% Espacamento de paragrafo
\setlength{\parindent}{0pt}
\setlength{\parskip}{7pt}

\widowpenalty=10000
\clubpenalty=10000
\flushbottom
